\section{Background}\label{background}

\subsection{Scalar comfort indices and what they compress}

PMV remains the most widely recognized steady-state comfort index for mechanically conditioned buildings \citep{Fanger1970PMV,ISO7730_2005,ASHRAE2023}. Adaptive comfort models provide another scalar framework by relating acceptable indoor temperatures to recent outdoor conditions \citep{deDear1998Adaptive,Brager2000Adaptive,Humphreys2002Adaptive}. These models are valuable for standards, simulation diagnostics, and building-performance comparison. Their strength is also their limitation: they compress a distributed occupant response into a scalar value, band, or violation state.

That compression is acceptable when the question is whether an average condition is near a standard reference. It is less complete when the question is whether severe sensation is probable. A scalar value cannot show whether probability mass is concentrated near neutral, spread across mild categories, or split between warm and cold discomfort tails. This paper does not argue that PMV or expected TSV should be discarded. It asks what additional information becomes visible when the full TSV probability distribution is retained.

\subsection{Probabilistic and ordinal TSV modeling}

The ASHRAE Global Thermal Comfort Database II provides a large field dataset linking environmental and personal variables to observed thermal sensation votes \citep{Foldvary2018GlobalDB}. Because TSV is ordered categorical rather than continuous, ordinal probabilistic models are a natural fit \citep{McCullagh1980Ordinal,Agresti2010Ordinal,Favero2023Ordinal}. They can preserve the order of the scale while producing class probabilities for each sensation category. Probability calibration is then needed because raw classifier scores are not automatically reliable probabilities \citep{Niculescu-Mizil2005PredictingGP,Gneiting2007Scoring}.

Recent comfort-modeling studies increasingly use probabilistic or Bayesian representations to express uncertainty in occupant response \citep{Luo2019GlobalDB,Kim2019Bayesian,Zhang2024BayesianMeta}. The contribution here is not a new claim that probabilistic comfort modeling exists. The contribution is a diagnostic use of the distribution: expected TSV and discomfort-tail probability are compared directly under future-weather stress, and the information loss from scalar and spatial aggregation is quantified.

\subsection{Future-weather stress tests and fixed mappings}

Climate-informed building simulation often uses generated or selected future weather files to test how current building assumptions respond under altered outdoor forcing \citep{GuoHe2026sAMY,GuoHe2026InputQuality}. Overheating and climate-responsive envelope studies similarly use future scenarios to expose vulnerability, while recognizing that actual future risk will depend on adaptation, envelope changes, operating practices, and occupant behavior \citep{Chaer2026OverheatingFramework,Ozarisoy2026ClimateResponsiveEnvelope}.

This study follows that stress-test logic. It does not assume that 2080 occupants have the same expectations, clothing behavior, physiology, or adaptive opportunities as current field observations. It holds the comfort-risk mapping fixed so the diagnostic question remains controlled. The result is conditional: under selected future-weather files, with a fixed prototype and fixed occupant model, how often do scalar summaries and mean aggregation hide modeled discomfort-tail exceedance?

\subsection{From diagnostic states to downstream uses}

Probabilistic comfort states could be used by many downstream decision frameworks, including rule-based supervisory policies, stochastic MPC, fuzzy rule bases, or operator dashboards. This paper does not compare those frameworks and does not claim stability, optimality, or actuator feasibility. Its scope ends at the diagnostic interface: the probability distribution, the scalar and tail summaries derived from it, and the audits that show where information is lost.
